% SPDX-License-Identifier: Apache-2.0
% covers: Fb
\datasheet{Fb}{Razboj's framebuffer: a memory of one word per pixel behind an AXI peripheral end.}

\dsfacts{\code{gpu/razboj/src/fb.rs}}{\code{razboj::fb::Fb}}%
{\code{//docs:razboj}, \code{//docs:axi}}

\subsection*{Function}

\code{Fb} is \code{N} words of memory, \code{px}, behind the
peripheral end of an AXI link. A pixel is one word. The rasteriser
writes pixels into it, and in Razboj's runs the same memory also
holds the display list and its count, which the rasteriser reads.

It is the peripheral client of an \code{AxiPer}, and its link is one
port, \code{bus}, a \code{PerPort} of the four channels: requests
come on \code{bus\_req} and write beats on \code{bus\_w}, and answers
go out on \code{bus\_ans} and \code{bus\_r}. It answers write
bursts and read bursts of any length. The rasteriser writes a run of a
row's pixels as one write burst (issue 987) and fetches a display list
entry as one read burst. A read's first beat is answered from the
memory in the cycle the request is taken and the rest follow a beat a
cycle, one burst at a time. A write is held while its beats arrive,
since AXI4 puts no identifier on the write data channel, and is
answered once, with its last beat. A beat writes only the bytes its
strobes name, so a beat with none writes nothing. It counts what it
served: \code{rbursts} read bursts taken, \code{rbeats} read beats
sent, \code{wbursts} write bursts taken and \code{wbeats} write beats
that wrote something, which a run reads to say how much of its time
the link was busy and how many pixels it wrote.

\subsection*{Parameters}

\begin{center}\footnotesize
\begin{tabular}{@{}lp{0.7\textwidth}@{}}
\toprule
Parameter & Meaning \\
\midrule
\code{A} & The width of the link's address, in bits. \\
\code{I} & The width of the AXI identifier, in bits. \\
\code{N} & Words of memory, a power of two. \\
\bottomrule
\end{tabular}
\end{center}

The tables use \code{Fb<16, 2, 1024>}, the framebuffer of
\code{//gpu/razboj:trace}.

\subsection*{Address map}

The word a burst names is its address shifted down by two, cut to
sixteen bits, and masked with \code{N}~$-$~1. The memory itself gives
no word a meaning. In \code{//gpu/razboj:trace} and the tests of
\code{gpu/razboj/src/sim.rs}, with \code{N} of 1024:

\begin{center}\footnotesize
\begin{tabular}{@{}lp{0.6\textwidth}@{}}
\toprule
Byte address & What \code{sim.rs} puts there \\
\midrule
\code{0x000} & The framebuffer, sixteen by sixteen pixels, a word
each. \\
\code{0x400} & The display list, eight words an instruction. \\
\code{0x600} & The count of instructions, written last. \\
\bottomrule
\end{tabular}
\end{center}

\dstables{Fb}

\subsection*{Behaviour}

\begin{itemize}
\item A read is taken when \code{bus\_r} has room and no write is waiting
for its beat. It is answered on \code{bus\_r} in the cycle it is taken,
with the word from \code{px}, \code{OKAY} and \code{last} set.
\item A write is taken when no write is held. Its first word and its
identifier go into \code{paddr} and \code{pid}, \code{pend} is set,
and \code{wbursts} counts it.
\item While a write is held, each of its beats is taken as it arrives,
and the last only when \code{bus\_ans} has room. A beat writes the
bytes its strobes name over the word at \code{paddr}, which then
moves to the next word; \code{wbeats} counts a beat with any strobe
set.
\item The last beat ends the write: \code{pend} clears and
\code{bus\_ans} sends \code{OKAY} with \code{pid}.
\item While a write is held, no request of either kind is taken.
\item An address beyond the memory wraps rather than ending the run.
\item \code{pixel} reads a word back for a run or a test.
\end{itemize}

\subsection*{Verification}

\begin{itemize}
\item \code{//gpu/razboj:razboj\_test} renders every scene on the rasteriser
through the link into \code{Fb}, reads the pixels back and compares
each with the model in \code{gpu/razboj/src/model.rs}. The tests are
\code{a\_scene\_of\_every\_kind\_agrees\_with\_the\_model},
\code{a\_clear\_says\_only\_its\_colour},
\code{a\_triangle\_is\_the\_same\_whichever\_way\_it\_is\_wound},
\code{a\_triangle\_hanging\_off}\allowbreak\code{\_the\_screen\_is\_clipped},
\code{a\_single\_pixel\_and\_an\_empty\_box} and
\code{pseudorandom\_scenes\_agree\_with\_the\_model}.
\item The build simulates \code{Fb<16, 2, 1024>}'s netlist, with the
display list and its count in \code{px}, against the trace of
\code{//gpu/razboj:trace} under nvc and Verilator:
\code{//docs:razboj\_sim\_fb\_fb\_tb\_test} and
\code{//docs:razboj\_vsim\_fb\_test}.
\item \code{//gpu/razboj:trace} and \code{//gpu/razboj:demo} assert that the
framebuffer equals the model's picture. The build runs
\code{//gpu/razboj:demo} to draw the picture in \code{//docs:razboj}.
\end{itemize}

\subsection*{Used by}

\code{razboj::sim::run} and \code{razboj::sim::run\_list}, and so Razboj's
tests, \code{//gpu/razboj:trace} and \code{//gpu/razboj:demo}. The system's
\code{//soc:demo} draws the core's display list a second time with
\code{run\_list}, on a plain link into \code{Fb}, and checks that the
picture is the same. On the lattice itself the memory is the
simulation-only \code{Ram}.

\subsection*{Limits}

It holds one write at a time and takes no read while it does, and a
read burst's beats are answered a cycle apart. A write burst's beats
are taken as they come; its length is not read, and its last beat ends
it. The
word index is sixteen bits wide, so the framebuffer reaches at most
65536 words.
